\uuid{8xrq}
\exo7id{4510}
\titre{exo7 4510}
\auteur{quercia}
\organisation{exo7}
\datecreate{2010-03-14}
\isIndication{false}
\isCorrection{false}
\chapitre{Suite et série de fonctions}
\sousChapitre{Convergence simple, uniforme, normale}
\cpge{mpsi-mp2i.limites-continuite.c, pcsi.limites-continuite.c, ptsi.limites-continuite.c, mpsi-mp2i.integration.c, pcsi.integration.b, ptsi.integration.b, bcpst1.limites-continuite.c, bcpst1.integration-segment.a, mp-mpi.suites-series-fonctions.a, pc.suites-series-fonctions.a, psi.suites-series-fonctions.a, pc.integration.a, psi.integration.a, bcpst2.integrales-generalisees}
\module{Analyse}
\niveau{L2}
\difficulte{}

\contenu{
\texte{
Soit $f : {\R^+} \to \R$ continue, non identiquement nulle, telle que $f(0) = 0$
et $f(x) \to 0$ lorsque $x\to+\infty$.

On pose $f_n(x) = f(nx)$ et $g_n(x) = f\left(\frac xn\right)$.
}
\begin{enumerate}
    \item \question{Donner un exemple de fonction $f$.}
    \item \question{Montrer que $f_n$ et $g_n$ convergent simplement vers la fonction nulle,
    et que la convergence n'est pas uniforme sur $\R^+$.}
    \item \question{Si $ \int_{t=0}^{+\infty} f(t)\,d t$ converge, chercher
       $\lim_{n\to\infty}  \int_{t=0}^{+\infty} f_n(t)\,d t$ et
       $\lim_{n\to\infty}  \int_{t=0}^{+\infty} g_n(t)\,d t$.}
\end{enumerate}
}
